%-------------------------------------------------------------------------------
%   APPENDIX A - ХАВСРАЛТ
%-------------------------------------------------------------------------------

\section*{Хавсралт А: Нэмэлт материалууд}
\addcontentsline{toc}{section}{Хавсралт А: Нэмэлт материалууд}

\subsection*{А.1 Flashcard dataset-ийн жишээ}

Дараах хүснэгтэд Egüü системд ашигласан flashcard dataset-ийн жишээг харуулав:

\begin{longtable}{|l|l|l|p{4cm}|l|}
\caption{Flashcard dataset-ийн жишээ} \\
\hline
\textbf{ID} & \textbf{Монгол бичиг} & \textbf{Фонетик} & \textbf{Англи орчуулга} & \textbf{Категори} \\
\hline
\endfirsthead

\multicolumn{5}{c}%
{{\tablename\ \thetable{} -- үргэлжлэл}} \\
\hline
\textbf{ID} & \textbf{Монгол бичиг} & \textbf{Фонетик} & \textbf{Англи орчуулга} & \textbf{Категори} \\
\hline
\endhead

\hline
\endfoot

001 & ᠠ & a & a (vowel) & vowels \\
\hline
002 & ᠡ & e & e (vowel) & vowels \\
\hline
003 & ᠢ & i & i (vowel) & vowels \\
\hline
004 & ᠣ & o & o (vowel) & vowels \\
\hline
005 & ᠤ & u & u (vowel) & vowels \\
\hline
006 & ᠥ & ö & ö (vowel) & vowels \\
\hline
007 & ᠦ & ü & ü (vowel) & vowels \\
\hline
008 & ᠪ & b & b (consonant) & consonants \\
\hline
009 & ᠴ & ch & ch (consonant) & consonants \\
\hline
010 & ᠳ & d & d (consonant) & consonants \\
\hline
011 & ᠭ & g & g (consonant) & consonants \\
\hline
012 & ᠬ & kh & kh (consonant) & consonants \\
\hline
013 & ᠯ & l & l (consonant) & consonants \\
\hline
014 & ᠮ & m & m (consonant) & consonants \\
\hline
015 & ᠨ & n & n (consonant) & consonants \\
\hline
\end{longtable}

\subsection*{А.2 Firebase Security Rules}

Egüü системд ашигласан Firebase Security Rules-ийн бүрэн код:

\begin{lstlisting}[language=JavaScript, caption=Firestore Security Rules]
rules_version = '2';
service cloud.firestore {
  match /databases/{database}/documents {
    
    // Flashcards collection - бүх хүн унших боломжтой
    match /flashcards/{flashcardId} {
      allow read: if true;
      allow write: if request.auth != null && 
                      request.auth.token.admin == true;
    }
    
    // Users collection - зөвхөн өөрийн өгөгдөл
    match /users/{userId} {
      allow read, write: if request.auth != null && 
                            request.auth.uid == userId;
    }
    
    // UserProgress collection - зөвхөн өөрийн прогресс
    match /userProgress/{userId} {
      allow read, write: if request.auth != null && 
                            request.auth.uid == userId;
    }
    
    // Favorites collection - зөвхөн өөрийн дуртай жагсаалт
    match /favorites/{userId} {
      allow read, write: if request.auth != null && 
                            request.auth.uid == userId;
    }
    
    // DailyWords collection - бүх хүн унших боломжтой
    match /dailyWords/{date} {
      allow read: if true;
      allow write: if request.auth != null && 
                      request.auth.token.admin == true;
    }
  }
}
\end{lstlisting}

\subsection*{А.3 Хэрэглэгчийн санал асуулгын загвар}

Туршилтын үед ашигласан санал асуулгын асуултууд:

\begin{enumerate}
    \item \textbf{Ерөнхий мэдээлэл:}
    \begin{itemize}
        \item Нас: \_\_\_\_
        \item Мэргэжил: \_\_\_\_
        \item Монгол бичиг мэдэх түвшин: Эхлэгч / Дунд / Ахисан
    \end{itemize}
    
    \item \textbf{Апп-ийн үнэлгээ (1-5 оноо):}
    \begin{itemize}
        \item Ерөнхий сэтгэл ханамж: \_\_\_\_
        \item Хэрэглэхэд хялбар байдал: \_\_\_\_
        \item UI/UX дизайн: \_\_\_\_
        \item Функцийн бүрэн байдал: \_\_\_\_
        \item Гүйцэтгэл (хурд): \_\_\_\_
        \item OCR таних нарийвчлал: \_\_\_\_
    \end{itemize}
    
    \item \textbf{Нээлттэй асуултууд:}
    \begin{itemize}
        \item Апп-ийн хамгийн таалагдсан тал юу вэ?
        \item Апп-ийн хамгийн таалагдаагүй тал юу вэ?
        \item Ямар функц нэмэхийг хүсч байна вэ?
        \item Энэ апп-ийг найз нөхөддөө санал болгох уу? (Тийм/Үгүй)
        \item Нэмэлт санал, сэтгэгдэл:
    \end{itemize}
\end{enumerate}

\subsection*{А.4 Системийн архитектурын дэлгэрэнгүй диаграмм}

\begin{figure}[H]
    \centering
    \begin{tikzpicture}[
        node distance=1.5cm,
        box/.style={rectangle, draw=black, fill=blue!20, minimum width=2.5cm, minimum height=0.8cm, align=center, font=\small},
        db/.style={cylinder, draw=black, fill=green!20, minimum width=1.8cm, minimum height=1.2cm, align=center, shape border rotate=90, font=\small},
        cloud/.style={ellipse, draw=black, fill=yellow!20, minimum width=2.5cm, minimum height=1.2cm, align=center, font=\small}
    ]
    
    % Client Layer
    \node[box] (ui) {UI Layer\\(React Native)};
    \node[box, below=0.8cm of ui] (context) {State Management\\(Context API)};
    \node[box, below=0.8cm of context] (services) {Services Layer};
    
    % Services
    \node[box, below left=1cm and 1cm of services] (auth-service) {Auth\\Service};
    \node[box, below=1cm of services] (firestore-service) {Firestore\\Service};
    \node[box, below right=1cm and 1cm of services] (ocr-service) {OCR\\Service};
    
    % Firebase
    \node[cloud, below=2cm of firestore-service] (firebase) {Firebase\\Platform};
    \node[db, below left=1cm and 0.5cm of firebase] (firestore-db) {Firestore};
    \node[box, below=1cm of firebase] (auth-fb) {Authentication};
    \node[db, below right=1cm and 0.5cm of firebase] (storage) {Storage};
    
    % TensorFlow
    \node[box, right=3cm of ocr-service] (tf-model) {TensorFlow.js\\Model};
    
    % Connections
    \draw[->, thick] (ui) -- (context);
    \draw[->, thick] (context) -- (services);
    \draw[->, thick] (services) -- (auth-service);
    \draw[->, thick] (services) -- (firestore-service);
    \draw[->, thick] (services) -- (ocr-service);
    
    \draw[->, thick] (auth-service) -- (firebase);
    \draw[->, thick] (firestore-service) -- (firebase);
    \draw[->, thick] (ocr-service) -- (tf-model);
    
    \draw[->, thick] (firebase) -- (firestore-db);
    \draw[->, thick] (firebase) -- (auth-fb);
    \draw[->, thick] (firebase) -- (storage);
    
    \end{tikzpicture}
    \caption{Системийн дэлгэрэнгүй архитектур}
    \label{fig:detailed-architecture}
\end{figure}

\subsection*{А.5 Төслийн статистик}

\begin{table}[H]
\centering
\caption{Төслийн хөгжүүлэлтийн статистик}
\begin{tabular}{|l|r|}
\hline
\textbf{Шалгуур} & \textbf{Утга} \\
\hline
Нийт кодын мөр & 15,420 \\
\hline
TypeScript файл & 87 \\
\hline
React компонент & 42 \\
\hline
Service файл & 12 \\
\hline
Utility функц & 28 \\
\hline
Test файл & 15 \\
\hline
Flashcard тоо & 520 \\
\hline
Категори тоо & 7 \\
\hline
Дэмжигдсэн хэл & 2 (Монгол, Англи) \\
\hline
Хөгжүүлэлтийн хугацаа & 3 сар \\
\hline
Туршилтын хэрэглэгч & 20 \\
\hline
\end{tabular}
\end{table}

\subsection*{А.6 Ашигласан технологиудын бүрэн жагсаалт}

\begin{table}[H]
\centering
\caption{Dependencies жагсаалт}
\begin{tabular}{|l|l|}
\hline
\textbf{Package} & \textbf{Хувилбар} \\
\hline
react-native & 0.81.5 \\
\hline
expo & \textasciitilde54.0.20 \\
\hline
firebase & 12.4.0 \\
\hline
@tensorflow/tfjs & 4.20.0 \\
\hline
@tensorflow/tfjs-react-native & 1.0.0 \\
\hline
@react-navigation/native & 7.1.8 \\
\hline
@react-navigation/native-stack & 7.1.9 \\
\hline
@react-navigation/bottom-tabs & 7.4.0 \\
\hline
nativewind & 4.2.1 \\
\hline
react-i18next & 15.2.0 \\
\hline
@expo/vector-icons & 15.0.3 \\
\hline
expo-image & 3.0.10 \\
\hline
expo-camera & 17.0.8 \\
\hline
expo-font & 14.0.9 \\
\hline
expo-haptics & 15.0.7 \\
\hline
\end{tabular}
\end{table}

\subsection*{А.7 Төслийн GitHub repository}

Төслийн эх код нь дараах хаягаар нээлттэй эхтэй байршуулагдсан:

\texttt{https://github.com/username/eguu-mongolian-script}

Репозиторид дараах материалууд багтсан:
\begin{itemize}
    \item Бүрэн эх код
    \item README.md (суулгах заавар)
    \item Flashcard dataset (JSON формат)
    \item TensorFlow OCR загварын код
    \item Туршилтын үр дүнгийн мэдээлэл
    \item Хэрэглэгчийн гарын авлага
\end{itemize}

\vfill
